% ---------------------------------------------------------------------------
%  Copyright / colophon
% ---------------------------------------------------------------------------
\thispagestyle{empty}
\vspace*{\fill}
\noindent
{\small
\textit{Architettura dei calcolatori: prima programmare, poi formalizzare.}\\
Libro di testo delle lezioni di ARC12, Architettura dei Calcolatori e
Laboratorio di Architettura dei Calcolatori, Università degli Studi di Napoli
«Parthenope», anno accademico 2026/27. Edizione italiana di \emph{Computer
Architecture: Program First, Formalise Later}.\par
\vspace{1em}

\noindent
Il libro segue, capitolo per capitolo, le venticinque presentazioni del corso:
il Capitolo~$N$ è la Lezione~$N$. Le presentazioni, i programmi di riferimento
del progetto \vermone e i test che assemblano ed eseguono ogni listato sono
conservati insieme ai sorgenti di questo libro; i listati stampati qui sono
letti dagli stessi file che i test eseguono. Per questo i listati, con i loro
commenti, sono in inglese in entrambe le edizioni: sono il codice provato, non
una sua copia.\par
\vspace{1em}

\noindent
Il libro è costruito attorno a \prg, un runtime didattico aperto per giochi in
assembly \riscv e in C, che ha come destinazioni il microcontrollore
ESP32\nobreakdash-C6, l'emulatore QEMU di Espressif e l'host desktop \prgqt. La
piattaforma, il suo codice sorgente e la sua documentazione sono distribuiti
con licenza MIT dai loro autori: Raffaele Montella (responsabile scientifico e
del progetto), Ivan Cafiero e Simone Boscaglia (studenti contributori),
Università degli Studi di Napoli «Parthenope». A chi usa \prg in esercitazioni,
relazioni, tesi o materiale didattico si chiede di citarlo attraverso i
metadati \fname{CITATION.cff} del suo repository \citep{prg32repo}. Il volume
che lo accompagna, \emph{Playing with RISC-V} \citep{prg32textbook}, è il
manuale di laboratorio della piattaforma.\par
\vspace{1em}

\noindent
L'implementazione di riferimento di \vermone, il progetto che accompagna
tutto il libro (i quattro programmi dell'Appendice~\ref{app:verm1}), è un
contributo di Giuseppe Coviello.\par
\vspace{1em}

\noindent
Il testo e le figure originali di questo libro sono rilasciati con licenza
Creative Commons Attribuzione -- Condividi allo stesso modo 4.0 Internazionale
(CC BY-SA 4.0). I programmi sono rilasciati con licenza MIT. Le fotografie
storiche sono riprodotte dalle slide delle lezioni a scopo didattico e restano
di proprietà dei rispettivi titolari.\par
\vspace{1em}

\noindent
RISC\nobreakdash-V è un marchio registrato di RISC\nobreakdash-V
International. ESP32, ESP32\nobreakdash-C3, ESP32\nobreakdash-C6 ed
ESP\nobreakdash-IDF sono prodotti e marchi di Espressif Systems. GitHub è un
marchio di GitHub, Inc. Tutti gli altri marchi appartengono ai rispettivi
titolari e sono usati qui al solo scopo di identificazione.\par
\vspace{1em}

\noindent
Edizione compilata il \today. Composto con \LaTeX{} in Latin Modern.\par
}
\vspace*{\fill}
\clearpage
